\section{Agreement}\label{sec:agreement}
%% FUNCTION: Shows that round robin achieves EFX less often in a world where agents agree more on
%% what is valuable, first in the reference cell, then as a count over all cells. The only proven
%% case, full agreement without divisibility, is deferred to Section 4.4. Survival curves show
%% that failures also become more severe. Because the size of the effect depends on how far
%% values are spread, the next section turns to dispersion.

Agreement describes how similarly the agents rank the goods. Heirs who all want the same
painting agree strongly; flatmates with different tastes agree little. In our model, agreement
is the parameter $\rho$, which we report as the Spearman rank correlation of the world from which
the instances are drawn (Section~\ref{sec:instances}). We vary it over its five levels, Spearman
$0$, $0.24$, $0.48$, $0.73$ and $1$, and keep everything else at the reference cell: $n=4$
agents, $m=12$ goods and Gini $0.52$ ($\theta=1$).

\subsection{EFX Rate}\label{sec:agreement-rate}

As Figure~\ref{fig:agreement} shows, Round Robin achieves EFX less often in a world with more
agreement. Without agreement, $0.741$ of the instances are EFX. The EFX rate falls to $0.515$ at
Spearman $0.24$, to $0.293$ at Spearman $0.48$ and to $0.109$ at Spearman $0.73$. Under full
agreement it is $0.0003$, which are $6$ of $20{,}000$ instances. In absolute numbers, Round Robin
fails EFX in $5{,}178$, $9{,}706$, $14{,}135$, $17{,}829$ and $19{,}994$ of the $R=20{,}000$
instances per level. The decline is steep from the start: each of the first two steps costs about
$0.22$, and at Spearman $0.48$ the EFX rate is already less than half of its value without agreement. This result is
measured, not proven. Its direction is plausible, because under agreement the good that one agent
picks tends to be the good that the next agent wanted as well. Later agents then lose goods they
value highly to earlier agents, and exactly these goods make the earlier bundles attractive to
them. We state this explanation as a conjecture.

\begin{figure}[ht]
  \centering
  \includegraphics{fig41_agreement_rate.pdf}
  \caption{EFX rate of Round Robin against the Spearman rank correlation of the world from which
  the instances are drawn, for $n=4$ agents, $m=12$ goods and Gini $0.52$ ($\theta=1$). Each
  point rests on $R=20{,}000$ instances.}
  \label{fig:agreement}
\end{figure}

The decline is not a property of the reference cell alone. For each of the $42$ pairs $(n,m)$
and each of the six levels of dispersion, we follow the EFX rate over the five levels of
agreement. In $240$ of these $252$ courses, the EFX rate falls strictly at each of the four
steps (measured). The $12$ exceptions are exactly the courses with $q=0$ and Gini $0.06$
($\theta=0.1$). There the EFX rate equals $1$ at every level of agreement, so it cannot fall;
Section~\ref{sec:dispersion-limit} explains this ceiling. The $252$ courses are not independent
confirmations, since all cells share the same random seeds (Section~\ref{sec:design}).

Only one case of this pattern is proven. Under full agreement, Round Robin is almost surely not
EFX whenever the number of goods is not divisible by the number of agents
(Section~\ref{sec:divisibility-identical}). The reference cell is divisible, so this result does
not apply to it, and the six EFX instances at Spearman $1$ show that failure is not certain
there. For divisible cells under full agreement, the EFX rate depends on dispersion
(Section~\ref{sec:dispersion}).

\subsection{Severity of Failures}\label{sec:agreement-severity}

The EFX rate counts how often Round Robin fails, but not by how much. Figure~\ref{fig:agreement-survival}
therefore shows, for every level of agreement, the share of instances with $\astar\geq\alpha$.
Each curve ends at $\alpha=1$ in the EFX rate of Figure~\ref{fig:agreement}. The earlier a curve
drops below $1$ and the lower it runs, the more severe are the failures.

\begin{figure}[ht]
  \centering
  \includegraphics{fig41_agreement_survival.pdf}
  \caption{Share of instances with $\astar\geq\alpha$ under Round Robin, one curve per level of
  Spearman rank correlation, for $n=4$ agents, $m=12$ goods and Gini $0.52$ ($\theta=1$). Each
  curve rests on $R=20{,}000$ instances and ends at $\alpha=1$ in the EFX rate. The point on
  each curve marks the $5\%$ quantile of $\astar$ among the instances that are not EFX.}
  \label{fig:agreement-survival}
\end{figure}

The curves are ordered over the whole range: a curve for more agreement never lies above a curve
for less agreement by more than $0.0001$. For small $\alpha$ they lie close together, with
shares between $0.997$ and $0.928$ at $\alpha=0.2$, and from there on they fan out. At $\alpha=0.8$, the share of instances that reach at least this
level is $0.846$ without agreement and $0.667$, $0.457$, $0.243$ and $0.018$ at the four higher
levels. At $\alpha=0.6$, the shares are $0.925$, $0.812$, $0.651$, $0.468$ and $0.191$. More
agreement therefore does not only move instances from EFX to slightly below EFX. It moves them
into the middle range of $\astar$. The curve for full agreement has a different shape from the
others: it falls steeply between $\alpha=0.3$ and $\alpha=0.7$ and is almost flat above $0.8$.
In this world nearly every instance fails, and most failures are moderate. A share of $0.592$ of
the instances still reaches $\astar\geq0.4$, but only $0.018$ reaches $\astar\geq0.8$.

Very severe failures remain rare at every level. The share of instances with $\astar<0.2$ rises
from $0.003$ without agreement to $0.072$ under full agreement. This is why the conditional
$5\%$ quantile, the point on each curve, moves little: it is $0.32$, $0.27$, $0.23$, $0.21$ and
$0.18$ for the five levels. Read alone, it would suggest that failures hardly become more severe.
The curves show otherwise. The median of $\astar$ among the instances that are not EFX falls
from $0.74$ to $0.68$, $0.60$, $0.54$ and $0.44$. In summary, in a world with more agreement, Round
Robin fails more often, and it also fails by more (measured).

Two restrictions apply. The curves share the same random seeds, so their distances are estimated
precisely, but the curves are not independent samples. Moreover, we read the curves of one cell.
The same ordering appears in the neighbouring cell with $m=13$ goods and for Gini $0.28$
($\theta=0.5$), where all curves lie higher: for $m=12$, the median among the instances that are
not EFX then falls from $0.90$ to $0.77$, and at least $0.994$ of the instances reach
$\astar\geq0.4$ at every level of agreement. We did not examine other cells.

\paragraph{Relation to literature.}
The two ends of our agreement axis are the models \textit{full independence} and
\textit{full correlation} of \textcite[p.~74]{bouveret2011}. They observe that full correlation
is the worst case for the utility that agents obtain from a picking sequence, and they mention a
model with intermediate correlation without developing it. Our results cover this intermediate
range for a different target, EFX, and show that the EFX rate deteriorates gradually and not only
at the end point. \textcite[p.~11]{branzei2025} remarks that Round Robin fails to produce an EFX
allocation in many instances, especially for correlated valuations, without reporting numbers.
Our measurements quantify this remark for small instances. Both her correlated model
\parencite[p.~8]{branzei2025} and the CORRELATED model of \textcite[p.~1406]{dickerson2014}
combine a common value of a good with an individual term, as ours does. As for existence, the
reference cell with $n=4$ and $m=12$ lies in the region where the existence of an EFX allocation
is open for every level of agreement below Spearman $1$. There we assume existence
(Section~\ref{sec:existence}), and under this assumption every failure reported here is a failure
of Round Robin. Under full agreement, existence is proven
\parencite{plaut2020}, with an alternative proof for additive valuations by
\textcite[Prop.~1]{branzei2025}, so the $19{,}994$ failures at Spearman $1$ are failures of Round
Robin without any assumption.

\paragraph{Recommendation.}
The less the people involved agree on what is valuable, the more Round Robin can be relied on to
deliver EFX; where they agree strongly, Round Robin should be expected to miss EFX, and by a
clear margin. How large these effects are depends on how far apart the values of the goods lie,
which is the subject of the next section.
